\usepackage[a4paper,top=20mm,bottom=22mm,left=22mm,right=22mm,headheight=20pt,headsep=8pt]{geometry}
\usepackage[T1]{fontenc}
\usepackage{lmodern}
\usepackage[english]{babel}
\usepackage{microtype}
\usepackage{graphicx}
\usepackage{amsmath,amssymb,amsthm,amsfonts,mathtools}
\usepackage{braket}
\usepackage{booktabs,tabularx,array,multirow}
\usepackage{enumitem}
\usepackage{listings}
\usepackage{adjustbox}
\usepackage{caption}
\usepackage{subcaption}
\usepackage{xcolor}
\usepackage{tikz}
\usetikzlibrary{quantikz,patterns,patterns.meta,arrows.meta,positioning,fit,calc}
\usepackage{pgfplots}
\usepgfplotslibrary{groupplots}
\pgfplotsset{compat=1.18}
\usepackage{titlesec}
\usepackage{fancyhdr}
\usepackage{tocloft}
\usepackage[most]{tcolorbox}
\usepackage{csquotes}
\usepackage{hyperref}
\usepackage[capitalize,nameinlink,noabbrev]{cleveref}
\usepackage[backend=biber,style=numeric-comp,sorting=none,maxbibnames=99]{biblatex}
\addbibresource{references.bib}

% A restrained, colorblind-friendly palette used consistently throughout.
\definecolor{AccentBlue}{HTML}{2457C5}
\definecolor{DeepBlue}{HTML}{173B5E}
\definecolor{BaselineBlue}{HTML}{0072B2}
\definecolor{InterceptOrange}{HTML}{E69F00}
\definecolor{AdvancedTeal}{HTML}{009E73}
\definecolor{ProposalPurple}{HTML}{6F5AA8}
\definecolor{WarningRed}{HTML}{B33A3A}
\definecolor{SoftBlue}{HTML}{EEF4FA}
\definecolor{SoftOrange}{HTML}{FFF6E5}
\definecolor{SoftTeal}{HTML}{EAF7F2}
\definecolor{SoftGray}{HTML}{F4F5F7}
\definecolor{MidGray}{HTML}{6B7280}

\hypersetup{
  colorlinks=true,
  linkcolor=AccentBlue,
  citecolor=AdvancedTeal,
  urlcolor=AccentBlue,
  pdftitle={Known-State Qubits as Both Resource and Risk},
  pdfauthor={Vi Connelly; GPT-5.6 Thinking (OpenAI), AI research and drafting collaborator},
  pdfsubject={Quantum compilation, qubit reuse, information flow, and verification}
}

\setlength{\parindent}{1.25em}
\setlength{\parskip}{0pt}
\setlength{\emergencystretch}{2em}
\setlist{nosep,leftmargin=1.7em,topsep=3pt}
\captionsetup{font=small,labelfont=bf,labelsep=period,skip=5pt}
\setcounter{tocdepth}{2}
\setcounter{secnumdepth}{3}
\renewcommand{\cftsecfont}{\color{AccentBlue}}
\renewcommand{\cftsecpagefont}{\color{AccentBlue}}
\renewcommand{\cftsubsecfont}{\small}
\renewcommand{\cftsubsecpagefont}{\small}
\setlength{\cftbeforesecskip}{2pt}
\setlength{\cftbeforesubsecskip}{0pt}

\titleformat{\section}
  {\Large\bfseries\color{DeepBlue}}
  {\thesection}{0.65em}{}
  [\vspace{0.2em}\color{black!25}\titlerule]
\titleformat{name=\section,numberless}
  {\Large\bfseries\color{DeepBlue}}
  {}{0pt}{}
  [\vspace{0.2em}\color{black!25}\titlerule]
\titleformat{\subsection}
  {\large\bfseries}
  {\thesubsection}{0.55em}{}
\titleformat{\subsubsection}
  {\normalsize\bfseries\color{DeepBlue}}
  {\thesubsubsection}{0.5em}{}
\titleformat{\paragraph}[runin]
  {\normalsize\bfseries}
  {}{0pt}{}[\quad]
\titlespacing*{\section}{0pt}{1.25em}{0.65em}
\titlespacing*{\subsection}{0pt}{1.0em}{0.4em}
\titlespacing*{\subsubsection}{0pt}{0.8em}{0.25em}
\titlespacing*{\paragraph}{0pt}{0.65em}{0.5em}

\pagestyle{fancy}
\fancyhf{}
\fancyhead[L]{\small\color{AccentBlue}Known-State Qubits as Both Resource and Risk}
\fancyhead[R]{\small\color{AccentBlue}$\triangleright$\ \nouppercase{\rightmark}}
\fancyfoot[C]{\thepage}
\renewcommand{\headrulewidth}{0.45pt}
\renewcommand{\headrule}{\hbox to\headwidth{\color{black!35}\leaders\hrule height \headrulewidth\hfill}}
\renewcommand{\sectionmark}[1]{\markright{\thesection\ #1}}
\fancypagestyle{plain}{
  \fancyhf{}
  \fancyfoot[C]{\thepage}
  \renewcommand{\headrulewidth}{0pt}
}

\newcommand{\Tr}{\operatorname{Tr}}
\newcommand{\id}{\mathbb{I}}
\newcommand{\E}{\mathbb{E}}
\newcommand{\Prb}{\operatorname{Pr}}
\newcommand{\norm}[1]{\left\lVert #1\right\rVert}
\newcommand{\ketproj}[1]{\ket{#1}\!\bra{#1}}
\newcommand{\cnot}{\operatorname{CNOT}}
\newcommand{\cz}{\operatorname{CZ}}
\newcommand{\todo}[1]{\textcolor{WarningRed}{\textbf{[TODO: #1]}}}

\theoremstyle{plain}
\newtheorem{theorem}{Theorem}[section]
\newtheorem{proposition}[theorem]{Proposition}
\newtheorem{corollary}[theorem]{Corollary}
\theoremstyle{definition}
\newtheorem{definition}[theorem]{Definition}
\theoremstyle{remark}
\newtheorem{remx}[theorem]{Remark}
\newenvironment{rem}
  {\pushQED{\qed}\renewcommand{\qedsymbol}{$\diamond$}\remx}
  {\popQED\endremx}

\newtcolorbox{plainbox}[1][]{
  enhanced,breakable,colback=SoftGray,colframe=black!45,
  boxrule=0.45pt,arc=1.5pt,left=7pt,right=7pt,top=5pt,bottom=5pt,
  before skip=5pt,after skip=5pt,title={#1},fonttitle=\bfseries
}
\newtcolorbox{claimbox}[1][]{
  enhanced,breakable,colback=SoftBlue,colframe=BaselineBlue,
  boxrule=0.55pt,leftrule=2.4pt,arc=1.5pt,left=7pt,right=7pt,top=5pt,bottom=5pt,
  before skip=5pt,after skip=5pt,title={#1},fonttitle=\bfseries
}
\newtcolorbox{cautionbox}[1][]{
  enhanced,breakable,colback=SoftOrange,colframe=InterceptOrange,
  boxrule=0.55pt,leftrule=2.4pt,arc=1.5pt,left=7pt,right=7pt,top=5pt,bottom=5pt,
  before skip=5pt,after skip=5pt,title={#1},fonttitle=\bfseries
}
\newtcolorbox{limitbox}[1][]{
  enhanced,breakable,colback=SoftGray,colframe=WarningRed,
  boxrule=0.55pt,leftrule=2.4pt,arc=1.5pt,left=7pt,right=7pt,top=5pt,bottom=5pt,
  before skip=5pt,after skip=5pt,title={#1},fonttitle=\bfseries
}
\newtcolorbox{aibox}[1][]{
  enhanced,breakable,colback=SoftTeal,colframe=AdvancedTeal,
  boxrule=0.55pt,leftrule=2.4pt,arc=1.5pt,left=7pt,right=7pt,top=5pt,bottom=5pt,
  before skip=5pt,after skip=5pt,title={#1},fonttitle=\bfseries
}

\newcommand{\evidencebadge}[2]{%
  \tikz[baseline=(n.base)]\node[rounded corners=1.5pt,inner xsep=4pt,inner ysep=1.5pt,
  fill=#1!13,draw=#1!75,text=#1!80!black,font=\scriptsize\bfseries] (n) {#2};%
}

\lstdefinestyle{python}{
  language=Python,
  basicstyle=\ttfamily\footnotesize,
  keywordstyle=\color{AccentBlue}\bfseries,
  commentstyle=\color{AdvancedTeal!75!black}\itshape,
  stringstyle=\color{ProposalPurple!85!black},
  showstringspaces=false,
  breaklines=true,
  frame=single,
  rulecolor=\color{black!25},
  backgroundcolor=\color{SoftGray},
  columns=fullflexible,
  xleftmargin=0.4em,
  xrightmargin=0.4em
}

\graphicspath{{./figures/}}
\numberwithin{equation}{section}
\numberwithin{figure}{section}
\numberwithin{table}{section}
